\section{CONCLUSIONES}

\USFQNotaPlantilla{En esta sección el autor debe responder a la(s) pregunta(s) de investigación con base en la información recabada, combinando la revisión de literatura con los resultados obtenidos en la sección de análisis. Se deben señalar las limitaciones del estudio realizado: número de personas, rango de edad, institución, actualización de la información, alcance geográfico o cualquier otra restricción pertinente.}

\USFQNotaPlantilla{Se debe describir la importancia de los datos y análisis del estudio, explicar a la audiencia por qué el estudio es significativo, quién se beneficiará del estudio, cómo cambiarán las opiniones acerca del campo de estudio que existían anteriormente y cuáles presunciones se cambiarán o confirmarán con el estudio. También se deberán explicar las posibles limitaciones metodológicas de diseño identificadas. Se deben incluir recomendaciones para futuros estudios. Esta sección puede estar dividida en subtítulos.}

\subsection{Respuesta a la pregunta de investigación}

\USFQCampo{Responder de manera directa la pregunta o preguntas de investigación.}

\subsection{Limitaciones}

\USFQCampo{Declarar las limitaciones del estudio y explicar cómo afectan la interpretación de los resultados.}

\subsection{Importancia y recomendaciones}

\USFQCampo{Explicar la importancia del estudio, sus beneficiarios potenciales y las recomendaciones para investigaciones futuras.}
